\chapter{Wyniki eksperymentów}

Rozdział przedstawia wyniki benchmarku pełnoskalowego oraz analizę końcowej serii
HPO. Wszystkie rankingi jakości pochodzą z jednego protokołu v5.3 i zostały
wyznaczone za pomocą oficjalnych metryk Numerai. Starszych uruchomień kontrolnych
nie włączono do tej populacji.

\section{Weryfikacja pełnej skali}
Benchmark pełnoskalowy objął 2\,746\,268 wierszy uczących, 780 cech i 574 ery.
Potok z dziesięcioma partycjami XGBoost wykonał trening, predykcję dla
3\,895\,460 wierszy walidacyjnych, ocenę i zapis diagnostyki w 1888,6~s
(31,5~min). Utworzone zostały konfiguracja rozwiązana, hash konfiguracji i zapis
przebiegu. Rezultat dokumentuje działanie tego samego frameworku na pełnym zbiorze,
nie jest porównaniem jakości ani identycznej konfiguracji z końcową serią HPO.

\begin{table}[htbp]
\centering
\caption[Wynik benchmarku pełnoskalowego]{Wynik operacyjny benchmarku pełnoskalowego}
\label{tab:full-scale-result}
\begin{tabular}{lr}
\toprule
Wielkość & Wartość \\
\midrule
Wiersze train / validation & 2\,746\,268 / 3\,895\,460 \\
Liczba cech / er train & 780 / 574 \\
Modele składowe & 10 $\times$ XGBoost \\
Czas całkowity & 1888,6 s (31,5 min) \\
\bottomrule
\end{tabular}
\sourceown
\end{table}

\section{Kompletność końcowej serii HPO}
Proces przerwano po około 1~h 43~min, po uruchomieniu \FinalTrialsLaunched{} prób,
a więc przed planowanym limitem 200 prób lub 11,5~h. Kompletny wynik uzyskano dla
\FinalTrialsCompleted{} konfiguracji (\FinalCompletionRate\% całej puli) i tylko
one weszły do dalszej analizy. Ich łączny czas modelowy wyniósł
\FinalCompletedComputeHours{}~h, a mediana czasu próby \FinalMedianRuntime{}~s.
Nierówne liczebności rodzin ograniczają porównanie do charakteru opisowego.

\section{Rozkład oficjalnych metryk}
Tabela~\ref{tab:metric-summary} syntetyzuje wszystkie kompletne konfiguracje.
Na holdoucie dodatni CORR Sharpe uzyskało 25 z 27 prób. Na późniejszym przedziale
selekcyjnym dodatni CORR Sharpe osiągnęły 24 próby, natomiast dodatni MMC Sharpe
13 prób. Różnica pokazuje, że sama korelacja z celem nie gwarantuje unikalności
sygnału względem Meta Modelu.

\begin{table}[htbp]
\centering
\caption[Podsumowanie oficjalnych metryk]{Podsumowanie oficjalnych metryk dla 27 kompletnych prób}
\label{tab:metric-summary}
\begin{tabular}{lrrr}
\toprule
Metryka & Średnia & Mediana & Wynik dodatni \\
\midrule
CORR Sharpe, holdout (52 ery) & 0,479 & 0,553 & 25/27 \\
CORR Sharpe, późniejszy (88 er) & 0,128 & 0,128 & 24/27 \\
MMC Sharpe, późniejszy (88 er) & 0,0247 & $-0,0098$ & 13/27 \\
\bottomrule
\end{tabular}
\sourceown
\end{table}

\section{Rodziny modeli w ustalonym budżecie}
Nazwa rodziny oznacza bazowy estymator użyty wewnątrz
\path{EraEnsembleModel}, a nie pojedynczy model. Wśród kompletnych prób
CatBoost wystąpił 10 razy (dziewięć zespołów po 5 modeli bazowych i jeden po 3),
LightGBM 6 razy (cztery zespoły po 5 i dwa po 3), a XGBoost 8 razy (trzy zespoły
po 5 i pięć po 3). Każdy taki zespół zawierał ponadto meta-model Ridge.
TabPFN wystąpił dwukrotnie, a TabICL raz.

Najwyższe mediany CORR Sharpe na holdoucie uzyskały CatBoost (0,659, $n=10$)
i LightGBM (0,655, $n=6$). Mediana XGBoost wyniosła 0,273 ($n=8$), TabICL
0,300 ($n=1$), a TabPFN 0,015 ($n=2$). Wszystkie obserwacje i rozrzut przedstawia
rysunek~\ref{fig:model-family-comparison}.

\begin{figure}[htbp]
\centering
\caption[Wyniki według rodziny modelu]{Oficjalny CORR Sharpe na holdoucie według
rodziny modelu. Punkty oznaczają kompletne próby, a liczności podano nad grupami.}
\label{fig:model-family-comparison}
\includegraphics[width=\textwidth]{model_family_comparison}
\sourceown
\end{figure}

\FloatBarrier
Na późniejszym przedziale mediana LightGBM wyniosła 0,280, CatBoost 0,133, XGBoost 0,036,
TabPFN 0,108, a TabICL 0,088. Wynik jest opisem badanego budżetu, nie rankingiem
architektur w ogólności: TPE przydzielił rodzinom różne liczby prób, a modele typu
foundation korzystały z limitów liczby przykładów i redukcji wymiaru.

\section{Przebieg optymalizacji}
Faza inicjalna trwała do uzyskania 20 kompletnych konfiguracji. Ostatnią
propozycją startową była próba 22, a pierwszą adaptacyjną propozycją TPE próba 23.
Lider holdoutu pojawił się w próbie 11.
Jeszcze w fazie startowej próba 22 podniosła maksimum CORR Sharpe na późniejszym
przedziale z 0,402 do 0,425. Siedem kompletnych prób adaptacyjnych
(23-27, 29 i 30) nie poprawiło żadnego z tych maksimów
(rysunek~\ref{fig:hpo-progression-final}).

\begin{figure}[htbp]
\centering
\caption[Przebieg końcowej serii HPO]{Wyniki prób i maksimum narastające dla
holdoutu oraz późniejszego przedziału. Linia przerywana oddziela fazę inicjalną
od adaptacji TPE.}
\label{fig:hpo-progression-final}
\includegraphics[width=\textwidth]{official_hpo_progression}
\sourceown
\end{figure}

\FloatBarrier
W zrealizowanym fragmencie fazy adaptacyjnej nie pojawił się nowy lider.
Obserwacja ta nie dowodzi nieskuteczności TPE: faza obejmowała tylko siedem
kompletnych prób, a nie wykonano równoległego eksperymentu kontrolnego z losowym
próbkowaniem o tym samym budżecie.

\section{Scenariuszowa ekstrapolacja dłuższego HPO}
Na prośbę o ocenę możliwego efektu dłuższego uruchomienia wykonano symulację
scenariuszową wyłącznie dla maksimum CORR Sharpe na historycznym holdoucie.
Nie jest to prognoza jakości rynkowej. Wariant empiryczny losuje z siedmiu
zaobserwowanych wyników adaptacyjnego TPE i nie tworzy niewidzianego ogona.
Wariant ostrożny estymuje średnią i odchylenie przez bootstrap tych samych siedmiu
wyników, a następnie wygładza rozkład rozkładem normalnym. Wariant optymistyczny
stosuje tę samą procedurę do 16 obserwacji LightGBM i CatBoost, czyli rodzin, które
w zrealizowanej serii miały najwyższe mediany. Dla każdego budżetu wykonano
200\,000 replikacji z ziarnem 20260830. W każdej replikacji wynik końcowy jest
większą z wartości obecnego lidera 1,0443 i maksimum przyszłych losowań.

\begin{figure}[htbp]
\centering
\caption[Ekstrapolacja budżetu HPO]{Scenariuszowa ekstrapolacja najlepszego CORR
Sharpe na holdoucie i odsetka replikacji poprawiających obecnego lidera. Pasma oznaczają
kwantyle 5-95\% między replikacjami, a nie przedziały ufności jakości live.}
\label{fig:hpo-budget-extrapolation}
\includegraphics[width=\textwidth]{hpo_budget_extrapolation}
\sourceown
\end{figure}

\begin{table}[htbp]
\centering
\footnotesize
\caption[Wyniki ekstrapolacji HPO]{Scenariusze maksimum CORR Sharpe na holdoucie
po zwiększeniu łącznej liczby kompletnych prób}
\label{tab:hpo-budget-extrapolation}
\begin{tabular}{llrrrr}
\toprule
Scenariusz & Budżet & Średnia & Mediana & Kwantyle 5-95\% & Replikacje z poprawą \\
\midrule
Ostrożny TPE & 50 & 1,0481 & 1,0443 & 1,0443-1,0676 & 8,0\% \\
Ostrożny TPE & 100 & 1,0542 & 1,0443 & 1,0443-1,1151 & 19,9\% \\
Ostrożny TPE & 200 & 1,0635 & 1,0443 & 1,0443-1,1488 & 33,5\% \\
Optymistyczny - drzewa & 50 & 1,0634 & 1,0443 & 1,0443-1,1700 & 24,1\% \\
Optymistyczny - drzewa & 100 & 1,0886 & 1,0443 & 1,0443-1,2398 & 47,0\% \\
Optymistyczny - drzewa & 200 & 1,1156 & 1,0837 & 1,0443-1,2889 & 63,5\% \\
\bottomrule
\end{tabular}
\sourceown
\end{table}

Empiryczne przelosowanie siedmiu wyników TPE daje dla każdego budżetu maksimum
1,0443 i zerowy odsetek replikacji z poprawą, ponieważ najlepszy zaobserwowany wynik tej fazy
jest niższy od dotychczasowego lidera. Dopiero założenie o ciągłym, niewidzianym
ogonie rozkładu dopuszcza poprawę. Duża różnica między scenariuszami oraz szerokie
górne granice pokazują, że wynik jest bardzo wrażliwy na założenia przy zaledwie
siedmiu próbach adaptacyjnych. Nie ekstrapolowano późniejszego CORR ani MMC,
ponieważ nie były funkcją celu HPO i służyły już do selekcji kandydatów.

\section{Transfer między przedziałami czasowymi}
Korelacja Pearsona między CORR Sharpe na holdoucie i późniejszym przedziale wyniosła
$r=0{,}360$, a korelacja rang Spearmana $\rho=0{,}377$. Związek jest dodatni,
lecz umiarkowany: wysoka pozycja na 52-erowym holdoucie tylko częściowo współwystępuje
z dobrym wynikiem później i nie zachowuje rankingu konfiguracji
(rysunek~\ref{fig:holdout-validation-final}).

\begin{figure}[htbp]
\centering
\caption[Holdout a późniejszy przedział]{Zgodność CORR Sharpe między holdoutem
i późniejszym przedziałem dla 27 kompletnych konfiguracji.}
\label{fig:holdout-validation-final}
\includegraphics[width=0.92\textwidth]{holdout_validation_scatter}
\sourceown
\end{figure}

\FloatBarrier
Próba 11 maksymalizuje wynik na holdoucie, lecz na późniejszym przedziale ustępuje
próbom 22, 4 i 10. Uzasadnia to zastosowanie dwustopniowej selekcji w tym badaniu:
holdout steruje HPO, natomiast późniejszy przedział dostarcza dodatkowej informacji
o zachowaniu konfiguracji w kolejnym okresie. Ponieważ wykorzystano go do wyboru
kandydatów, nie pełni on funkcji niezależnego zbioru testowego.

\section{Kandydaci końcowi}
Tabela~\ref{tab:final-candidates} zestawia pięć reprezentatywnych konfiguracji.
Próba 11 jest liderem holdoutu, próba 22 liderem CORR na późniejszym przedziale, a próba 4
liderem MMC. Próby 10 i 8 pokazują mocne alternatywy z rodzin LightGBM i XGBoost.

\begin{table}[htbp]
\centering
\small
\caption[Końcowi kandydaci]{Oficjalne metryki końcowych kandydatów. HO - holdout,
PP - późniejszy przedział}
\label{tab:final-candidates}
\begin{tabular}{rlrrrrrr}
\toprule
Próba & Model & \multicolumn{2}{c}{HO CORR} & \multicolumn{2}{c}{PP CORR} & \multicolumn{2}{c}{PP MMC} \\
\cmidrule(lr){3-4}\cmidrule(lr){5-6}\cmidrule(lr){7-8}
& & średnia & Sharpe & średnia & Sharpe & średnia & Sharpe \\
\midrule
22 & LightGBM & 0,03032 & 0,7149 & \textbf{0,01679} & \textbf{0,4247} & 0,00802 & 0,2467 \\
4 & CatBoost & 0,02161 & 0,4779 & 0,01483 & 0,4024 & \textbf{0,01031} & \textbf{0,3063} \\
10 & LightGBM & 0,02743 & 0,7191 & 0,01450 & 0,3372 & 0,00716 & 0,1871 \\
11 & LightGBM & \textbf{0,03742} & \textbf{1,0443} & 0,01038 & 0,2880 & 0,00310 & 0,1023 \\
8 & XGBoost & 0,02047 & 0,5072 & 0,00705 & 0,1971 & 0,00142 & 0,0450 \\
\bottomrule
\end{tabular}
\sourceown
\end{table}

W przestrzeni CORR-MMC późniejszego przedziału tylko próby 22 i 4 nie są zdominowane
(rysunek~\ref{fig:validation-pareto-final}). Próba 22 maksymalizuje jakość sygnału,
a próba 4 jego unikalność. Nie istnieje więc jeden zwycięzca niezależny od celu
wdrożeniowego.

\begin{figure}[htbp]
\centering
\caption[Front Pareto CORR-MMC]{Front Pareto na późniejszym przedziale. Okręgi
wyróżniają próby niezdominowane względem CORR Sharpe i MMC Sharpe.}
\label{fig:validation-pareto-final}
\includegraphics[width=0.80\textwidth]{validation_pareto}
\sourceown
\end{figure}

\FloatBarrier
Do treningu nastawionego na stabilny CORR wybrano próbę 22. Jest to zespół pięciu
modeli bazowych LightGBM i meta-modelu Ridge,
z RobustScalerem i grupową selekcją wariancji. Router wskazuje 989 unikalnych
cech z grup
\texttt{medium+agility+midnight}, a po selekcji do LightGBM trafia 876 kolumn.
Model wykorzystuje 200 rund, tempo uczenia 0,0200, 31 liści i neutralizację cech
0,2373. Jej CORR Sharpe wynosi 0,7149 na
holdoucie i 0,4247 na późniejszym przedziale, przy dodatnim MMC Sharpe 0,2467.

Próba 4 jest alternatywą nastawioną na MMC. Wykorzystuje pięć modeli bazowych
CatBoost i meta-model Ridge, StandardScaler, 970 cech z grup
\texttt{small}, \texttt{agility}, \texttt{constitution}, \texttt{dexterity},
\texttt{faith} oraz \texttt{intelligence}, 2000 rund,
głębokość 5, tempo uczenia 0,0314
i neutralizację 0,4947. Osiąga CORR Sharpe 0,4024 oraz MMC Sharpe 0,3063
na późniejszym przedziale. Próba 11 uzyskała najwyższy wynik na holdoucie:
CORR Sharpe \FinalBestCorrSharpe{}, średni CORR \FinalBestMeanCorr{}, maksymalne
obsunięcie \FinalBestDrawdown{} i \FinalBestPositiveEras\% dodatnich er.
